\begin{lemma}[Penalization converges to the strict constrained maximum]
Let $Z$ be a metric space, let $K\subseteq Z$ be compact, and let
$\Theta,\chi:Z\to\R$ be continuous on $K$ with $\chi\ge0$ on $K$. Let $k_0\in K$
satisfy $\chi(k_0)=0$, and assume that $k_0$ is a \emph{strict} maximum point of
$\Theta$ among the points of $K$ at which $\chi$ vanishes:
\[
  \Theta(k)<\Theta(k_0)\qquad
  \bigl(k\in K,\ \chi(k)=0,\ k\ne k_0\bigr).
\]
Let $(\lambda_m)_{m\ge0}$ be positive reals with $\lambda_m\to\infty$, and let
$(k_m)_{m\ge0}$ be points of $K$ with
\[
  \Theta(k)-\lambda_m\chi(k)\le\Theta(k_m)-\lambda_m\chi(k_m)
  \qquad(k\in K,\ m\ge0),
\]
that is, $k_m$ maximizes $\Theta-\lambda_m\chi$ over $K$. Then $k_m\to k_0$.
\end{lemma}
\begin{proof}
\emph{Step 1: two consequences of maximality, taking $k=k_0$.}
The function $\Theta$ is continuous on the nonempty compact set $K$, so it
attains a maximum there (preamble), say at $k_\star\in K$; thus
$\Theta(k)\le\Theta(k_\star)$ for every $k\in K$. Taking $k=k_0$ in the
maximality hypothesis and using $\chi(k_0)=0$ gives
\begin{equation}\label{eq:pmc-base}
  \Theta(k_0)\le\Theta(k_m)-\lambda_m\chi(k_m)\qquad(m\ge0).
\end{equation}
From \eqref{eq:pmc-base} we read off two facts. First, since $\lambda_m>0$ and
$\chi(k_m)\ge0$ (as $k_m\in K$), the subtracted term is nonnegative, so
\begin{equation}\label{eq:pmc-value}
  \Theta(k_0)\le\Theta(k_m)\qquad(m\ge0).
\end{equation}
Second, since $\Theta(k_m)\le\Theta(k_\star)$,
\begin{equation}\label{eq:pmc-penalty}
  \lambda_m\chi(k_m)\le\Theta(k_\star)-\Theta(k_0)=:C_0\qquad(m\ge0),
\end{equation}
and $C_0\ge0$ because $k_0\in K$ gives $\Theta(k_0)\le\Theta(k_\star)$.

\emph{Step 2: the penalty vanishes.}
Dividing \eqref{eq:pmc-penalty} by $\lambda_m>0$,
\[
  0\le\chi(k_m)\le\frac{C_0}{\lambda_m}\qquad(m\ge0),
\]
the left inequality because $k_m\in K$. Since $\lambda_m\to\infty$ and $C_0$ is
a fixed real number, $C_0/\lambda_m\to0$; by the squeeze principle,
\begin{equation}\label{eq:pmc-chi}
  \chi(k_m)\longrightarrow0\qquad(m\to\infty).
\end{equation}

\emph{Step 3: every subsequence has a sub-subsequence converging to $k_0$.}
It suffices to prove this: a sequence in a metric space converges to $k_0$ as
soon as every subsequence admits a further subsequence converging to $k_0$.

So let $(n_j)_j$ be indices with $n_j\to\infty$. The points $k_{n_j}$ all lie in
the compact set $K$, so by sequential compactness of a compact metric space
(preamble) there are a strictly increasing $(m_j)_j$ and a point $\hat k\in K$
with $k_{n_{m_j}}\to\hat k$. We claim $\hat k=k_0$.

Since every $k_{n_{m_j}}$ lies in $K$ and $k_{n_{m_j}}\to\hat k$ with $\hat k\in
K$, the convergence also holds within $K$; as $\chi$ and $\Theta$ are continuous
on $K$, this gives
\[
  \chi\bigl(k_{n_{m_j}}\bigr)\longrightarrow\chi(\hat k),
  \qquad
  \Theta\bigl(k_{n_{m_j}}\bigr)\longrightarrow\Theta(\hat k).
\]
On the other hand $(n_{m_j})_j\to\infty$, because $n_j\to\infty$ and
$(m_j)_j$ is strictly increasing, so \eqref{eq:pmc-chi} applies along this
subsequence and gives $\chi(k_{n_{m_j}})\to0$. Limits in a metric space are
unique, hence
\[
  \chi(\hat k)=0 .
\]
Likewise \eqref{eq:pmc-value} holds for every index, so passing to the limit in
$\Theta(k_0)\le\Theta(k_{n_{m_j}})$ gives
\[
  \Theta(k_0)\le\Theta(\hat k).
\]
Now suppose $\hat k\ne k_0$. Then $\hat k\in K$ with $\chi(\hat k)=0$ and
$\hat k\ne k_0$, so the strict-maximum hypothesis gives
$\Theta(\hat k)<\Theta(k_0)$, contradicting the previous display. Therefore
$\hat k=k_0$, and $k_{n_{m_j}}\to k_0$.

As $(n_j)_j$ was an arbitrary sequence of indices tending to infinity, the
criterion quoted at the start of this step applies and gives $k_m\to k_0$.
\end{proof}
